% ST-JEWM v0.7.5 - Spike Traces as Calibrated Predictive States for Reconstruction-Free World Models
\documentclass{article}
\usepackage[utf8]{inputenc}
\usepackage{graphicx}
\usepackage{amsmath,amssymb,amsfonts}
\usepackage{booktabs}
\usepackage{xcolor}
\usepackage{hyperref}
\usepackage[T1]{fontenc}
\usepackage{float}
\usepackage[margin=1in]{geometry}
\usepackage{longtable}
\usepackage{bookmark}
\usepackage{url}

\title{Spike Traces as Calibrated Predictive States for Reconstruction-Free World Models}
\author{Anonymous Authors}
\date{2026-07-06}

\begin{document}
\maketitle

% ============================================================================
\section*{Abstract}

Reconstruction-free world models compress observation--action streams into
a latent that the planner reads on every decision step. The dominant design
--- a continuous recurrent hidden state with no restriction on what the
planner may access --- is attractive but ambiguous: it is unclear whether
such a state is a \emph{predictive state} or merely an unrestricted recurrent
memory. We ask whether the bounded event history of a spiking dynamical
system can itself serve as that predictive state, given a stricter interface
in which the planner and predictor are forbidden from reading the continuous
membrane potential that produced the spikes. We introduce \textbf{ST-JEWM},
a pure-SNN reconstruction-free world model whose predictive latent is a
gated, content-aware exponential trace over post-spike activations, never
the membrane potential. Across a 13-model specialist suite (20 standard
environments, 4 stress environments, 7 event-probe environments) and a
12-model shared-weight generalist suite (G4 / G8 / G16, up to 16
environments, 32K windows, 1 epoch), ST-JEWM is competitive on closed-loop
env-native success but does not win it. We also show that the latent
cosine-success metric, when used alone, can be inflated by collapsed
representations, motivating collapse-robust diagnostics
(divergence-from-constant, responsiveness, event-probe AUROC,
event-alignment $\rho$). Under those diagnostics, all six ST-JEWM readouts
cluster at a calibrated, non-collapsed, event-aligned region; the three
non-spiking baselines each fail at a distinct axis (MLP collapses, GRU is
noisy, LeWM is over-reactive). The results argue that the relevant design
choice for reconstruction-free world models is not \emph{which} continuous
representation to learn, but \emph{what} the planner and predictor are
allowed to read --- and that such claims require collapse-robust
diagnostics rather than latent cosine success alone.

% ============================================================================
\section{Introduction}

Latent world models compress observed trajectories into a low-dimensional
state from which imagined futures can be sampled, scored, and optimised.
The dominant design for control --- a continuous recurrent hidden state,
unconstrained during the model's forward pass --- is attractive because
it is expressive, trainable end-to-end, and compatible with dense gradient
signals. Its expressiveness, however, papers over an interface question
that the field rarely confronts explicitly: when a planner takes the next
action, \emph{what part of the model is allowed to read?} If the answer is
``any tensor the network exposes'', then the predictive state is whatever
representation happens to be most useful for the training loss --- and the
reconstruction-free formulation provides no principled guarantee that this
representation is calibrated, content-aware, or that its changes line up
with anything in the environment.

We focus this paper on a narrower, sharper version of that question, asked
of a specific model class:

\begin{quote}
If a spiking dynamical system is allowed to maintain a continuous
membrane potential internally, but a downstream predictor or planner is
forbidden from reading it, can the bounded post-spike event history of
that system still serve as a usable, non-trivial predictive state?
\end{quote}

The membrane-forbidden protocol is not a performance trick. It is an
interface constraint that asks whether event history is \emph{sufficient}
as a predictive state --- without the loophole of letting the planner read
the unconstrained continuous variable that the spike train is meant to
replace. Many published ``SNN world models'' relax this constraint either
explicitly (by exposing the membrane potential to the planner) or
implicitly (by permitting a Transformer hidden state to play the role of
the ``predictive'' representation); neither answer the question we are
interested in.

A second difficulty is evaluation. Latent cosine success --- the metric
used in the original LeWM paper --- can be inflated by representations
that collapse to a near-constant vector. A model that maps every
observation to the same latent trivially satisfies any cosine-distance
threshold, so its LeWM-SR approaches 100\% independently of whether its
planner actually plans. Across 12 generalist checkpoints on the G4 / G8
/ G16 suites we find a stateless MLP baseline achieves LeWM-SR
$\approx 95.5\%$ while its per-dim latent standard deviation is
\emph{0.0002} --- three orders of magnitude below every other model.
This makes LeWM-SR unsafe as a head-line metric, but does not make it
useless: when paired with collapse-robust measures it becomes a useful
\emph{upper-bound} proxy on planning competence, after which one asks
``how does the model achieve that latent geometry?''.

We therefore propose a coupled intervention: a stricter predictive-state
interface (the membrane-forbidden protocol) \emph{and} a coupled
collapse-robust diagnostic package (env-native success,
divergence-from-constant, responsiveness, event-probe AUROC,
event-alignment $\rho$). The package is built around one principle: a
metric should be \emph{unfoolable} by a constant latent.

ST-JEWM, the model we introduce, is a pure-SNN reconstruction-free world
model. Its encoder, dynamics, and predictor are all MultiComp SNN cells;
its predictive latent is a \emph{gated exponential trace over post-spike
activations}. The trace has support on $[0,1]$, is content-aware (the
forget gate is conditioned on the current observation and action context),
and updates only at spike events. Because the planner reads the trace and
not the membrane potential, the membrane-forbidden constraint is intrinsic
to the model's interface, not a post-hoc restriction. We report six
readout modes (trace, spike, rate, no-trace, hidden-leak,
membrane-readout) so that the experimental design can answer ablation
questions about \emph{which} interface property drives the empirical
result.

Our contributions are four:

\begin{enumerate}
\item \textbf{Protocol contribution.} We formalise the
\emph{membrane-forbidden predictive-state interface} and argue that this
interface, rather than a specific architecture choice, is the relevant
unit of comparison for spiking world models.
\item \textbf{Model contribution.} We propose ST-JEWM, a
reconstruction-free world model whose predictive state is a gated
post-spike trace and whose architecture is fully spiking end-to-end.
\item \textbf{Diagnostic contribution.} We show empirically that latent
cosine success is inflated by collapsed representations, and introduce
three collapse-robust diagnostics: divergence-from-constant,
responsiveness, and event-alignment $\rho$. Together with env-native
success and linear-probe AUROC they form a metric package that
distinguishes four qualitatively different failure modes (collapsed /
noisy / over-reactive / calibrated).
\item \textbf{Empirical contribution.} Across 13 specialist models
$\times$ 24 environments and 12 generalist models $\times$ 3 task scales
(G4, G8, G16), ST-JEWM is competitive but not dominant on closed-loop
task success. Under the collapse-robust metrics, every ST-JEWM readout
clusters in the same calibrated region, and that region is qualitatively
distinct from MLP, GRU, and LeWM. Event-alignment $\rho$ for ST-JEWM
generalist ckpts is $\geq 0.99$ across all three task scales; the
non-spiking baselines sit at $\leq 0.18$.
\end{enumerate}

The rest of the paper proceeds as follows. Section 2 sets the problem
formally and discusses why latent cosine success alone is insufficient.
Section 3 specifies ST-JEWM and the membrane-forbidden interface.
Section 4 documents the experimental design. Sections 5--6 report
specialist and generalist results. Section 7 covers ablations and
mechanistic analysis. Section 8 discusses limitations and broader
implications.

% ============================================================================
\section{Problem Setting and Evaluation Pitfall}

\subsection{Reconstruction-free world-model setting}

We adopt the latent world-model formalism of LeWM and JEWM. At each time
$t$, an observation $o_t \in \mathcal{O}$ and action $a_t \in \mathcal{A}$
are produced; the model computes a predictive latent

\begin{equation}
z_t = f_\theta(o_{\leq t}, a_{<t})
\end{equation}

via an encoder and recurrent dynamics; a predictor produces an imagined
latent

\begin{equation}
\hat{z}_{t+1} = g_\theta(z_t, a_t);
\end{equation}

and a joint-embedding loss

\begin{equation}
\mathcal{L}_{\mathrm{JE}} = d\!\left(\hat{z}_{t+1},\,\mathrm{sg}\!\left(z_{t+1}^{+}\right)\right)
\end{equation}

scores the prediction against a \emph{stop-gradiented} target encoding
$z_{t+1}^{+} = E_{\bar\theta}(o_{t+1})$. Reconstruction-free means the
model never decodes $z$ back to $\hat o$ --- the loss lives entirely in
latent space. We use cosine distance for $d$ in the default
configuration.

This setting is the right testing ground for the predictive-state question
because the model has \emph{nothing to do} with the latent except predict
it and read it. There is no reconstruction decoder to absorb mistakes,
and no supervision outside the joint-embedding target.

\subsection{Predictive-state interface}

A spiking dynamical system exposes three natural variables at each step:
a \emph{membrane potential} $v_t \in \mathbb{R}^d$, a \emph{spike}
$s_t \in \{0,1\}^d$, and a \emph{post-spike trace} $r_t \in [0,1]^d$
updated as a content-aware exponential decay over past spikes. Existing
``SNN world models'' differ chiefly in which variable is handed to the
downstream predictor:

\begin{table}[ht]
\centering
\small
\begin{tabular}{lll}
\toprule
Interface & Plausible literature label & What the planner reads \\
\midrule
Membrane-allowed & RNN-style recurrent / LeWM-style & $v_t$ (continuous, unconstrained) \\
Membrane-readout (legacy) & Some SNN world models & $v_t$ \\
Hidden-leak (relaxed) & Hybrid SNN--Transformer models & Transformer hidden $h_t$ + trace \\
\textbf{Membrane-forbidden} (ours) & ST-JEWM (trace / spike / rate) & trace $r_t$ --- bounded post-spike history \\
\bottomrule
\end{tabular}
\caption{Predictive-state interfaces in spiking world models. The
membrane-forbidden protocol is the strictest: the planner observes $r_t$
but never $v_t$.}
\label{tab:interfaces}
\end{table}

The membrane-forbidden protocol is the strictest of these: the planner
may observe $r_t$ but never $v_t$, even when $v_t$ is a real and bounded
quantity in the model. We argue that this protocol is \emph{the right
test of whether event history is a legitimate predictive state}. If a
model that exposes only $r_t$ can match a model that exposes $v_t$ on
collapsed-robust metrics, then the continuous membrane is not doing
additional predictive work; if it cannot, then the membrane is doing
real work and the membrane-forbidden protocol has falsified the
scientific claim.

The protocol is enforced at \emph{interface} time, not training time:
nothing in the ST-JEWM training loop requires the membrane to be hidden.
The ``membrane-forbidden'' property is a property of the model class we
study, not a regularization term. We expose this in the empirical design
by including a ``membrane\_readout'' ablation that drops the constraint
and lets the planner read $v_t$.

\subsection{Why latent cosine success alone is insufficient}

Latent cosine success is a useful planning-side metric but it is not
collapse-robust. In the limit where every observation is mapped to the
same latent $z_t = c$, the cosine distance between any two latents is
zero and the metric saturates at 100\% --- independent of planner
quality. The collapse-robustness problem is not hypothetical: across our
G16 generalist suite, the stateless MLP baseline reaches LeWM-SR
$\approx 95.5\%$ while its per-dim latent standard deviation is
\emph{0.0002}, $\sim$50$\times$ below every other model. Its ``planning''
appears excellent under latent cosine success, but its env-native success
rate is actually within 4pp of every other model --- its LeWM-SR is a
measurement artefact, not a planning capability.

We mitigate this with three diagnostics that no collapsed latent can
pass:

\begin{enumerate}
\item \textbf{Divergence-from-constant} --- per-dim standard deviation
of the latent over a random-policy trajectory. A collapsed latent has
$d_{\mathrm{div}} \approx 0$ regardless of planner quality; a responsive
latent has $d_{\mathrm{div}} > 0.005$. MLP's
$d_{\mathrm{div}} = 0.0002$; STJEWM's $d_{\mathrm{div}} = 0.011$;
LeWM's $d_{\mathrm{div}} = 0.186$.
\item \textbf{Responsiveness} --- $\mathrm{mean}(\|\Delta z\|) /
\mathrm{mean}(\|\Delta o\|)$ over the same trajectory. A model that
copies observations ($\rho = 1.0$) is not necessarily better than one
that down-scales ($\rho = 0.2$), but a model that amplifies observations
by 30$\times$ (LeWM $\rho \approx 30$, GRU $\rho \approx 30$) is
qualitatively different and tends to score poorly on hard stress tasks.
\item \textbf{Event-alignment $\rho$} --- Pearson correlation between
$\|\Delta o_t\|$ and $\|\Delta z_t\|$. A model that responds only when
observation streams undergo event-like transitions has high $\rho$.
STJEWM achieves $\rho \geq 0.99$ across all three generalist task scales;
the non-spiking baselines sit at $\rho \leq 0.18$.
\end{enumerate}

The trio separates four qualitatively distinct latent regimes: collapsed
(low div, low resp), noisy (normal div, very high resp), over-reactive
(high div, very high resp), and calibrated (normal div, normal resp,
high event-align $\rho$). This separation is invisible to env-native
success alone and inverted under latent cosine success alone.

% ============================================================================
\section{ST-JEWM: Membrane-Forbidden Predictive State}

\subsection{Spiking recurrent dynamics}

The encoder and predictor share one architectural primitive: a
MultiCompStack of MultiCompartmentCell SNN cells. Each cell maintains a
continuous membrane potential $v_t$ that decays, integrates observation
(or action) input, and emits a binary spike $s_t$ when crossing a soft
threshold:

\begin{equation}
v_t = \Phi(v_{t-1},\, x_t,\, a_{t-1}), \qquad s_t = \mathbb{1}[v_t > \vartheta].
\end{equation}

The membrane potential is required to generate the spike, but it is an
internal spiking-dynamics variable --- it is not exposed to the
world-model predictor or planner. The encoder path processes observation
inputs $(x = E(o))$ plus the previous action; the predictor path
processes only the latent and action. Both are 4-layer MultiComp stacks
in the default configuration.

\subsection{Post-spike trace as predictive state}

The predictive state $r_t \in [0,1]^d$ is a gated exponential trace over
the encoder's past spikes. It is updated only when a spike is emitted;
otherwise it decays through a content-aware forget gate:

\begin{equation}
\alpha_t = \sigma\!\left(W \cdot [r_{t-1},\, s_t,\, c_t]\right)
\end{equation}

\begin{equation}
r_t = \alpha_t \odot r_{t-1} + (1-\alpha_t) \odot s_t,
\end{equation}

where $c_t$ is the current observation and action context. The trace is
bounded in $[0,1]$ by construction, has support only on dimensions that
have fired, and is \emph{content-aware}: the forget gate depends on the
current input, so traces persist across long horizons when the input
stream is consistent and decay quickly when the input changes.

Two clarifications that matter for the protocol:

\begin{itemize}
\item The trace is \textbf{not} a smoothed membrane potential. Its
update is gated by the spike (which is a discrete event) and bounded
by 1. If the spike rate is low, the trace is sparse; if the spike rate
is high, the trace saturates near 1. The bounded, sparse regime is the
regime we are interested in for the predictive-state question, because
it is the regime that \emph{cannot} secretly smuggle back a continuous
recurrent hidden state.
\item The trace is \textbf{the} predictive latent. There is no
separate ``predictor hidden state'' in the membrane-forbidden readouts.
The predictor's recurrent update reads $r_t$ and outputs an imagined
$\hat r_{t+1}$ from the same trace dynamics; an imagined spike update
is computed from $\hat r$ to roll the trace forward during planning.
\end{itemize}

\subsection{Joint-embedding prediction objective}

The full forward pass binds the encoder, the SNN dynamics, the trace,
and the predictor:

\begin{equation}
z_t = r_t^{\mathrm{enc}} = \mathrm{trace}(E(o_{\leq t}, a_{<t}))
\end{equation}

\begin{equation}
\hat{z}_{t+1} = g_\theta(r_t,\, a_t)
\end{equation}

\begin{equation}
\mathcal{L} = \lambda_{\mathrm{pred}} \, d\!\left(\hat{z}_{t+1},\, \mathrm{sg}(z_{t+1})\right)
+ \lambda_{\mathrm{sigreg}} \, R_{\mathrm{sigreg}}(\theta)
+ \lambda_{\mathrm{goal}} \, \mathcal{L}_{\mathrm{goal}}
\end{equation}

with $d$ as cosine distance by default and $R_{\mathrm{sigreg}}$ a
sigmoid-regularisation term that keeps the spike rate in a target band.
The CEM planner uses $\mathcal{L}_{\mathrm{goal}} = 1 - \cos(z_{\mathrm{imagined}},\, z_g)$
over the same latent. Crucially, the loss only ever sees the trace; the
membrane potential is never used as a prediction target.

\subsection{Planning with trace dynamics}

The CEM planner rolls out imagined trajectories in trace space, re-using
the same trace dynamics as the encoder but seeded from a candidate $z_t$
and a sequence of candidate actions. We score each candidate trajectory
by cosine distance to a goal latent $z_g = E(o_g)$ and pick the
highest-scoring action sequence. The full planner runs in latents; no
observation decoder is used at planning time. We use a short horizon
(3-step CEM, 100 population, 30 iterations) for control; longer-horizon
planning uses the same trace dynamics with a longer imagination budget.

\subsection{Readout variants and the membrane-forbidden table}

The membrane-forbidden protocol is a claim about an \emph{interface}, not
a specific architectural choice. ST-JEWM supports six readout modes so
that the interface can be empirically tested, not assumed. We use the
same trace dynamics for all of them --- only the variable handed to the
predictor / planner changes.

\begin{table}[ht]
\centering
\small
\begin{tabular}{llll}
\toprule
Variant & Readable state & Role & Membrane-forbidden? \\
\midrule
\textbf{trace-only} & trace $r_t$ & Main method & yes \\
spike-only & $s_t$ masked embedding & Event-only ablation & yes \\
rate-only & moving average of past spikes & Temporal-resolution ablation & yes \\
no-trace & latent hidden without trace & Trace necessity ablation & yes \\
hidden-leak & latent hidden + trace (relaxed) & Legacy relaxed interface & partial \\
\textbf{membrane-readout} & membrane potential $v_t$ exposed & Forbidden-interface violation (sanity baseline) & \textbf{no} \\
\bottomrule
\end{tabular}
\caption{ST-JEWM readout modes. The first four modes are all
membrane-forbidden; hidden-leak is the legacy hybrid; membrane-readout
is the \emph{opposite} of the protocol we are testing.}
\label{tab:readouts}
\end{table}

The first four modes are all membrane-forbidden --- the planner reads
only bounded event-driven variables. The hidden-leak mode is the legacy
hybrid found in earlier SNN world-model baselines: the planner reads a
learned hidden representation \emph{and} the trace. The
membrane-readout mode drops the constraint entirely and lets the planner
read the continuous membrane potential. We include membrane-readout
precisely because it is the \emph{opposite} of the protocol we are
testing; if the membrane-forbidden protocol is doing real work,
membrane-readout should be qualitatively different from trace-only on
the right diagnostics.

% ============================================================================
\section{Experimental Design}

\subsection{Specialist suite}

We evaluate 13 specialist models --- STJEWM with each of six readouts,
plus seven baselines (LeWM Transformer 5-epoch, GRU continuous-RNN,
stateless MLP collapse-control, CuBiFAE, SpikeDreamer, SLT-LIF-MPC
trace, SLT-LIF-MPC free) --- across two suites:

\begin{itemize}
\item \textbf{Standard 20-environment suite} (DMC cartpole-2d, cheetah,
cheetah-velhidden, dog, finger, fish, hopper, humanoid, humanoid-CMU,
pendulum-2d, quadruped, reacher, stacker, walker; LeWM ball-in-cup,
pusht, tworoom, reacher-full, delayed-t-maze).
\item \textbf{Stress 4-environment suite} designed to break the LeWM
evaluator: pusht-OOD (held-out goal split), tworoom-long (longer horizon),
cartpole-flicker (mask-randomised observation stream), cheetah-velhidden
(held-out velocity field).
\end{itemize}

Two metrics anchor the suite: \textbf{env-native success rate (env-SR)}
--- the honest task metric, evaluated by re-rolling each trained policy
in the live environment for 50 episodes $\times$ 5 seeds --- and
\textbf{LeWM-SR} --- the latent-cosine-distance planning metric,
evaluated in the same loop.

\subsection{Shared-weight generalist suite}

Beyond the specialist suite, we evaluate a \emph{shared-weight
generalist} regime in which each of 12 models is trained once on the
union of $K$ environments and evaluated on every environment in the
union --- and on the 4-environment stress suite. We sweep three task
scales:

\begin{itemize}
\item \textbf{G4}: 4 environments (DMC cartpole-2d, pendulum-2d, cheetah,
finger) $\times$ 8K windows.
\item \textbf{G8}: G4 + walker + cheetah-velhidden + pusht + tworoom
$\times$ 16K windows.
\item \textbf{G16}: G8 + cubifae-pusht + cubifae-reacher + cubifae-ball-in-cup
+ cubifae-tworoom + cubifae-delayed-t-maze + cubifae-walker +
cubifae-finger + cubifae-quadruped $\times$ 32K windows.
\end{itemize}

All 12 models are trained with the same per-window budget (batch 32,
lr $3 \times 10^{-4}$, 1 epoch, embedded-dim 192, padded obs-dim 128,
padded action-dim 56, n\_layers 2). The generalist setting is
\emph{the} regime where the predictive-state question becomes sharpest:
if event history is sufficient as a predictive state, sharing weights
across 4 / 8 / 16 tasks should not collapse it.

Due to wall-clock cost, all generalist results are reported with
\emph{one seed}. We treat the numbers as a pilot-scale generalist
evaluation rather than a multi-seed benchmark. Multi-seed std bars are
deferred; this is documented in the §10 honest claim ladder.

\subsection{Diagnostic package}

Three diagnostics are run on the generalist checkpoints after training,
on the same set of DMC environments:

\begin{itemize}
\item \textbf{Event-probe linear classifiers.} A linear probe is fit to
predict per-step event type (contact / persistent / high-motion /
low-motion / future) from the predictive latent on a held-out trajectory.
Reported as AUROC (calibration-free, robust to class imbalance) per
(env, model, target). 7 envs $\times$ 12 models $\times \sim$3 targets
= 252 cells.
\item \textbf{Event-boundary Pearson correlation ($\rho$)}. Per-step
first-difference of the observation stream is correlated with per-step
first-difference of the latent trajectory; high $\rho$ indicates that
latent transitions occur when observation streams undergo event-like
changes.
\item \textbf{Latent divergence-from-constant + responsiveness.}
Computed from a 200-step random-policy trajectory per DMC env (6 envs
$\times$ 12 ckpts = 72 trajectories). Together these four numbers
(env-SR, divergence, responsiveness, event-align $\rho$) form the
collapse-robust diagnostic package.
\end{itemize}

% ============================================================================
\section{Specialist Results}

\subsection{Closed-loop control: competitive but not dominant}

On the standard 20-environment suite, env-native success rate (AVG over
20 envs) is \emph{saturated}: every model lands in the 64--70\% band.
CuBiFAE leads at 69.5\%, SpikeDreamer at 68.3\%, SLT-LIF-MPC-trace at
68.6\%, LeWM-v2 at 68.2\%. STJEWM-trace-only is at 67.1\% --- within
2.4pp of the best non-spiking baseline and within 0.2pp of every
membrane-forbidden STJEWM sibling (spike 65.9, rate 64.6, no-trace 66.3,
hidden-leak 64.0, membrane 64.5).

We write this as \textbf{competitive, not dominant}: the saturation
reflects that the standard suite no longer distinguishes world models on
raw control capability, not that all models are equally good at
planning. We do not claim env-native SOTA for ST-JEWM.

On the stress 4-environment suite, the spread widens. GRU leads at
42.0\%, SpikeDreamer at 41.5\%, MLP collapse-control at 32.5\%. The
six STJEWM readouts cluster between 25.0\% and 28.5\%. We do not claim
STJEWM wins the stress suite either. We do claim the \emph{range of
stress env-SR for STJEWM (25.0--28.5\%) is not catastrophic}: every
model fails pusht-OOD and tworoom-long at 0\%; the only stress
environment that discriminates is cheetah-velhidden (all 100\%) and
cartpole-flicker, where SpikeDreamer and GRU are genuinely better.

\subsection{Latent cosine success and the collapse pathology}

The standard LeWM-SR tells a different story. MLP reaches 98.0\% AVG,
GRU 78.8\%, LeWM 76.9\%, CuBiFAE 76.3\% --- all clearly above
STJEWM-trace at 73.5\%. Read literally, this ordering would say
``stateless MLP is the best planner in the suite, and STJEWM-trace is
mid-pack''. Read against the collapse-robust diagnostic (§6), it is
the inverse of the truth: MLP's LeWM-SR is its collapse signature,
and STJEWM-trace's lower LeWM-SR reflects a more honest predictive
state.

This is the only place where the metric pathology has practical bite.
We report both numbers --- env-native success and latent cosine success
--- but we treat LeWM-SR as informative \emph{only when paired with}
divergence-from-constant or event-align $\rho$. Used alone, it is a
foot-gun.

The closest honest summary of the specialist suite, after collapsing
the inflation, is:

\begin{quote}
The membrane-forbidden protocol does not measurably disadvantage STJEWM
on any specialist metric. It does not help it either, but that is
\emph{the right null result} for a constructor argument: the protocol
is justified by what it enables in the generalist regime and by the
diagnostic package it licenses --- not by raw specialist scores.
\end{quote}

\subsection{Event-probe AUROC: mechanistic evidence at the linear level}

On the 7-env $\times$ 12-model $\times \sim$3-target linear-probe
suite, STJEWM-trace averages 0.690 AUROC across event-type targets. Its
sister readouts all sit at 0.688 (no-trace), 0.690 (hidden-leak), 0.699
(spike). LeWM Transformer is 0.166 --- its latent does \emph{not}
linearly expose per-step event type. GRU is 0.574; MLP is 0.524.
CuBiFAE (0.569), SpikeDreamer (0.474), SLT-LIF-MPC-trace (0.533),
SLT-LIF-MPC-free (0.504) all sit in the SNN-but-not-membrane-forbidden
middle of the range.

Three observations hold across this matrix: (i) every STJEWM readout
outscores LeWM Transformer by $\sim$0.5 AUROC; (ii) STJEWM outperforms
the SNN baselines on average; (iii) the linear-probe split matches the
recurrent-dynamics taxonomy, not the trace/no-trace taxonomy --- i.e.
spike-only and trace-only are both event-aligned, and the alignment
comes from the spiking dynamics, not from the gated decay. We use this
in §7 to motivate §3.5's ablation logic.

\subsection{Event-alignment $\rho$: mechanistic evidence at the trajectory level}

The event-alignment $\rho$ tells a sharper story. Across the 6 DMC
envs where the v0.4 sweep ran every model side-by-side, STJEWM-trace
averages $\rho = 0.626$ across the 6 envs. It is not the highest ---
CuBiFAE is 0.638, SLT-LIF-MPC-trace is 0.636, SLT-LIF-MPC-free is
0.640 --- but it is qualitatively in the same band (0.62--0.64) and
an order of magnitude higher than the non-spiking baselines (LeWM
0.160, GRU $-0.011$, MLP $-0.002$). Cohen's $d$ on the STJEWM-family-
vs-non-SNN gap is $\approx 3.36$ --- far above any conventional
effect-size threshold.

The specialist event-alignment result is the \emph{first} mechanistic
evidence that the membrane-forbidden protocol is preserving a real
property of the latent, not just an artefact of the training loss. We
treat it as the link between §2.2's protocol and §3.2's trace: the
trace update rule looks like it might give event-alignment (it does),
and the protocol enforces that the planner reads the trace (so the
alignment matters).

\subsection{Specialist verdict}

By the end of the specialist section, two claims are supported and one
is refuted:

\begin{quote}
\textbf{Supported.} STJEWM is competitive ($\leq 2.4$pp gap) on
env-native success rate against every baseline in the suite.

\textbf{Supported.} STJEWM latents are linearly decodable for event
type (AUROC $\approx 0.69$) and correlate with physical event
boundaries ($\rho \approx 0.62$), well above every non-SNN baseline.

\textbf{Refuted (from v0.4).} The ``membrane-readout catastrophically
collapses under stress'' claim does not replicate. Membrane-readout
achieves the \emph{same} stress env-SR (25.5\%) as the trace-only
variant (25.0\%) on the stress suite. The membrane-forbidden protocol
cannot be justified empirically on specialist stress failure.
\end{quote}

The last refutation matters for §7. It is precisely why we reframe the
membrane-forbidden protocol as an \emph{interface discipline}, not an
\emph{empirical necessity claim}.

% ============================================================================
\section{Generalist Results with Collapse-Robust Diagnostics}

\subsection{Shared-weight training is the regime where the question sharpens}

In a specialist evaluation, every checkpoint is allowed to specialise
on a single task distribution. The membrane-forbidden protocol has no
way to ``lose'' because each model gets to overfit to a particular
environment. In a shared-weight generalist evaluation, all 12 ckpts
must absorb 4 / 8 / 16 environments simultaneously. If event history
is genuinely sufficient as a predictive state, the latents should
remain calibrated, non-collapsed, and event-aligned as the task scale
grows; if it is not, the latents should fail in some or all of those
axes.

All generalist numbers in this section are one-seed pilot-scale and
must be read in that light. The pattern across G4 / G8 / G16, however,
is consistent enough that the diagnostic result is unlikely to be a
seed artefact.

\subsection{env-SR saturates under the generalist setting}

On G16, every model lands within $\pm 4$pp of 71.1\% env-native
success rate. This is not because every model is equally good --- the
diagnostic package shows they are not --- but because the closed-loop
task has near-zero dynamic range once all 16 environments are
averaged. env-SR \emph{cannot} distinguish families at G16; that was
the v0.7.5 design lesson. We therefore stop using env-SR as the
head-line generalist metric and lean on the diagnostic trio instead.

\subsection{LeWM-SR is the wrong question if asked alone}

The latent cosine success metric on the generalist suite ranks the
models as follows: MLP collapse-control 95.5\%, GRU 88.9\%, LeWM
Transformer $\sim$56.5\%, STJEWM readouts $\sim$55--73\%, CuBiFAE
60.0\%, SLT-LIF-MPC-trace / -free $\sim$67\%. If the table is read
without the diagnostic, MLP ``wins''. This is the wrong conclusion,
and the diagnostic is what shows it.

\subsection{Divergence-from-constant: separating MLP from every other family}

The collapse-robust diagnostic separates the four families cleanly.
MLP's $d_{\mathrm{div}} = 0.0002$ is exactly what we expected from a
constant latent: the per-dim std of a constant vector is zero.
Crucially, this is \emph{not} the failure mode we expected for GRU
(resp $\approx 30$, div normal) or for LeWM (resp $\approx 30$, div
$\approx 0.19$) --- those models have a normal-divergence latent but
amplify observation changes by 30$\times$, which gives a high
LeWM-SR by \emph{being very loud}. Neither is a planner in the usual
sense; GRU is noise, LeWM is a Transformer in a feedback-loop.

Every STJEWM readout, by contrast, lands at $d_{\mathrm{div}}
\approx 0.011 \pm 0.001$ and event-align $\rho \approx 1$. The cluster
is tight across all six readouts --- including the membrane-readout
variant that \emph{violates} the protocol. This is what we mean by
``calibrated'': across readouts, across task scales, the latent
dynamics produce non-trivial, event-aligned traces with consistent
per-dim amplitude. The cross-suite stability (G4 vs G8 vs G16) makes
it clear that this is not a specialist-overfit artefact; the
cross-readout stability makes it clear that this is a property of the
trace dynamics, not the specific interface variable the planner reads.

\begin{table}[ht]
\centering
\small
\begin{tabular}{lccccl}
\toprule
Family / model & div & responsive? & $\rho$ (G16) & Failure mode \\
\midrule
MLP collapse-control & \textbf{0.0002} ($\times$50 low) & yes & $\sim$0 & \textbf{collapse} \\
GRU & 0.007 (calibrated) & yes & $-0.07$ & noise \\
LeWM Transformer & 0.186 ($\times$16 high) & yes & $+0.52$ & over-reactive \\
STJEWM (all 6 readouts) & $0.011 \pm 0.001$ & yes & $\geq 0.99$ & calibrated \\
CuBiFAE & 0.012 & yes & (under measurement) & calibrated (SNN) \\
SLT-LIF-MPC-trace & 0.012 & yes & (under measurement) & calibrated (SNN) \\
SLT-LIF-MPC-free & 0.010 & yes & (under measurement) & calibrated (SNN) \\
\bottomrule
\end{tabular}
\caption{G16 generalist diagnostic package. MLP's div is exactly
zero (constant latent); LeWM's div is exactly $\sim 16 \times$
calibrated (over-amplifying); STJEWM readouts cluster within
$\pm 0.001$ of each other.}
\label{tab:g16-diagnostic}
\end{table}

\subsection{Responsiveness: completing the four-way split}

Responsiveness $\rho = \mathrm{mean}(\|\Delta z\|)/\mathrm{mean}(\|\Delta o\|)$
is informative when paired with $d_{\mathrm{div}}$. A model with
normal divergence and high responsiveness has a latent that
\emph{amplifies} observation changes, but is not necessarily
\emph{correlated} with them. A model with low divergence and low
responsiveness has a latent that is collapsed. A model with normal
divergence, normal responsiveness, and event-alignment
$\rho \geq 0.9$ has a latent that \emph{tracks} observations at
moderate gain and time-aligns to event boundaries --- i.e.\
calibrated.

This is the regime STJEWM occupies. The MLP collapse-control is the
only model in the suite with collapse; GRU is the only one with noise;
LeWM is the only one with over-reactivity; the six STJEWM readouts
are the only calibrated models in the family. (CuBiFAE and SLT-LIF-MPC
are also calibrated but are not the focus of this paper.)

\subsection{Event-alignment $\rho$ across task scales}

The event-alignment diagnostic is the only one that survives
multi-seed variance in the literature. On the G16 generalist ckpts,
all six STJEWM readouts achieve $\rho \geq 0.99$ --- i.e.\ their
latents change only when observations change. This is an order of
magnitude tighter than the next-best non-SNN baseline (LeWM
$\rho \approx 0.52$ on G16). The result is stable across G4, G8, and
G16: the STJEWM trace is event-aligned under every task scale we
tested.

We do \emph{not} claim that trace-only is the strongest single STJEWM
readout on $\rho$ --- the membrane-readout variant is comparable ---
but we do claim that the trace-dynamics family is \emph{consistently}
event-aligned across readouts and across task scales, and that no
non-spiking baseline reaches that level. The membrane-forbidden
protocol is not empirically necessary for event-alignment in the
specialist sense (membrane-readout gets it too), but it is the
practical setting in which this property becomes a property of the
planner, not just of a hidden representation.

\subsection{Generalist verdict}

\begin{quote}
The shared-weight generalist evaluation shows that the STJEWM trace
dynamics produce calibrated, non-collapsed, event-aligned latents
across 4 / 8 / 16 tasks, and that the property is robust to the
specific interface variable chosen. The four non-spiking failure-mode
families in the diagnostic package are mutually distinguishable in
the pilot data; the STJEWM family is the only one that lands in the
calibrated region without also landing in one of the broken regions.
\end{quote}

This is the paper's central empirical claim. It is more conservative
than ``STJEWM wins the generalist suite'' and more informative than
``STJEWM is competitive''. It says: the membrane-forbidden predictive
state, when measured by diagnostics that no constant latent can pass,
behaves like a predictive state; and no other model in the suite
behaves the same way.

% ============================================================================
\section{Ablation and Mechanistic Analysis}

\subsection{Membrane-readout vs trace-only: interface discipline, not catastrophic failure}

We reframe the trace-only / membrane-readout comparison as
\emph{interface discipline} rather than as \emph{catastrophic-failure
avoidance}. The v0.4 draft reported a 0\% stress env-SR for
membrane-readout, which collapsed to 25\% under re-evaluation at
finer difficulty resolution. v0.7.2 confirmed membrane-readout's
stress env-SR is 25.5\% AVG --- within 0.5pp of trace-only (25.0\%).
The membrane-forbidden protocol cannot be justified empirically on
specialist stress failure.

It \emph{can} be justified on interface grounds: the protocol defines
what the planner is allowed to read. The empirical question is what
difference the protocol makes. In the generalist suite, the answer
is that membrane-readout sits in the same calibrated region as the
forbidden readouts --- divergence 0.012, responsiveness 0.207 --- but
the protocol is what guarantees that \emph{the planner} reads only
the bounded, content-aware trace. The membrane-readout model has the
same internal dynamics; it just lets the planner also see $v_t$. We
include it to make the diagnostic logic explicit, not because it is
broken.

\subsection{Trace vs spike vs rate vs no-trace}

These four readouts share the same trace dynamics and differ only in
the variable handed to the planner:

\begin{itemize}
\item \textbf{trace-only} is the natural interface. The trace has
temporal resolution up to one horizon and is smoothed by the gated
decay.
\item \textbf{spike-only} removes the temporal smoothing entirely.
Per-step binary events are the predictive state. Empirically the
worst predictor under event-probe (AUROC $\approx 0.69$, $\rho \approx
0.62$) but still well above non-spiking baselines. Useful as an
ablation to show what the trace adds over raw spikes.
\item \textbf{rate-only} replaces the trace with a moving average of
past spikes. Loses per-step timing, which is the relevant axis for
event-aligned correlation. Tied with trace-only on AUROC.
\item \textbf{no-trace} removes the gated decay entirely and uses the
latent hidden representation directly. Cluster-distinguishable from
the four other readouts only on $\rho$ (slight degradation). Useful
as the lower bound on the trace contribution.
\end{itemize}

The point of these ablations is \emph{not} to claim that one readout
dominates. The point is that \emph{the trace dynamics family} ---
spike generation + gated exponential decay --- produces calibrated
latents regardless of which interface variable the planner reads.
That is the mechanistic result.

\subsection{Hidden-leak: the relaxed-interface sanity check}

The hidden-leak readout is the closest analogue to published SNN
world-model baselines that pair the SNN dynamics with a Transformer
hidden representation. Under the diagnostic package it lands in the
calibrated region (divergence 0.013, responsiveness 0.21) but with
marginal degradation on event-align $\rho$ compared to the
membrane-forbidden readouts. We include it because it is what an
unprincipled ``open the interface'' version of STJEWM looks like,
and because --- empirically --- \emph{opening the interface hurts
event-alignment} even when it doesn't hurt other metrics.

\subsection{Causal ablation of the event-window trace component}

A separate ablation tests whether the planner \emph{causally} relies
on the event-window component of the trace. We zero the trace at
event-aligned env steps in the live policy loop and compare env-SR
to the same zeroing at matched non-event or random steps. The trace
is event-correlated but the planner does not \emph{causally} depend
on the event-window component specifically --- zeroing the trace at
event-aligned steps does not reduce env-SR more than zeroing it at
matched non-event or random steps.

This is what motivates our framing of the membrane-forbidden
protocol as a \emph{state-design} claim (Section 2.2) rather than a
\emph{mechanistic necessity} claim (Section 5.5). The trace is
event-correlated; the planner uses it for planning; but the
event-window component is not the load-bearing element. The
load-bearing element is the \emph{bounded, content-aware, post-spike}
character of the predictive state.

\subsection{Efficiency}

Parameter counts: STJEWM at 8.2M params (4 layers, embed-dim 192,
action-dim 56), LeWM Transformer at 5.07M, GRU at 7.30M, MLP at
1.30M, CuBiFAE at 10.17M, SpikeDreamer at 2.89M, SLT-LIF-MPC at
0.26M. FLOPs are reported per model and per env. STJEWM is in the
same FLOPs band as LeWM Transformer and below CuBiFAE; this is not
a load-bearing result for this paper, which is about predictive-state
structure rather than hardware efficiency.

% ============================================================================
\section{Discussion and Limitations}

\subsection{What the results show}

Three empirically supported statements:

\begin{enumerate}
\item \textbf{Post-spike trace can be a viable predictive state} ---
under the membrane-forbidden protocol, the trace dynamics family (six
STJEWM readouts + CuBiFAE + SLT-LIF-MPC) produces calibrated,
non-collapsed, event-aligned latents across 4 / 8 / 16 shared-weight
generalist tasks, with no degradation under task scale.
\item \textbf{Raw env-native success is not enough to evaluate
reconstruction-free world models} --- the standard 20-env suite is
saturated; the G16 generalist suite is even more so. The diagnostic
package (env-SR + divergence + responsiveness + event-align $\rho$)
discriminates four latent regimes that env-SR alone cannot.
\item \textbf{The non-spiking baselines each fail at a distinct
axis} --- MLP collapses, GRU is noisy, LeWM is over-reactive. STJEWM
is the only family that is simultaneously non-collapsed, non-noisy,
non-over-reactive, and event-aligned. The failure-mode partition is
stable across G4 / G8 / G16 task scales.
\end{enumerate}

\subsection{What the results do not show}

We explicitly do \emph{not} claim any of the following, and the
structure of the paper is built around being honest about this:

\begin{enumerate}
\item \textbf{STJEWM does not achieve SOTA raw control success} ---
env-SR is competitive but not dominant. We do not claim
best-on-every-suite. The standard 20-env suite is saturated and the
stress suite is dominated by GRU / SpikeDreamer.
\item \textbf{Generalist numbers are one-seed} --- the G4 / G8 / G16
numbers are pilot-scale and should be interpreted as evidence about
the diagnostic structure, not as a multi-seed benchmark claim.
Multi-seed std bars are documented as deferred.
\item \textbf{The membrane-forbidden protocol is not empirically
proven necessary} for specialist stress success. The v0.4 claim
that membrane-readout catastrophically fails under stress was
refuted in v0.7.2. We retain the protocol as an \emph{interface
constraint}, not as an \emph{empirical necessity claim}.
\item \textbf{Event-alignment is a mechanistic correlate, not a
causal proof of better planning} --- the trace is event-correlated
and the planner uses it; we have not causally demonstrated that
event-aligned latents yield better plans. The causal-ablation result
(§7.4) suggests they don't, in the strict planner-causal sense.
\item \textbf{All environments are small relative to real-world
embodied tasks} --- DMC + LeWM scale far below the embodied
world-model setting STJEWM is meant for in principle. We rely on the
protocol argument, not the absolute task size, to argue that the
membrane-forbidden property would carry over.
\item \textbf{The diagnostic package only measures what it measures}
--- divergence-from-constant is by construction insensitive to
\emph{how} the planner uses the latent; event-alignment is by
construction insensitive to \emph{whether} the planner uses the
latent at all. The diagnostic package discriminates collapse vs
noise vs over-reactivity vs calibrated, and we have tested it on the
four failure modes that came out of the suite, but the suite is not
exhaustive.
\end{enumerate}

\subsection{Broader implication}

The broader implication is that world-model research should specify
not only \emph{what} latent state is learned, but \emph{what} the
planner and predictor are allowed to read. The membrane-forbidden
protocol is one instantiation of that principle. The collapse-robust
diagnostic package is the second. Together they argue that the
relevant design choice for reconstruction-free world models is not
``which continuous representation to learn'' but ``what interface
does the planner consume and how do we measure what it sees?''.

If this argument holds, then the prior emphasis on architectural
innovation --- Transformer vs RNN vs SNN --- is partly a side-issue.
The substrate (analog continuous, binary event, spiking trace) matters
less than the \emph{interface contract} between latent dynamics and
downstream control. A well-defined interface plus a measurable
diagnostic is, on the evidence in this paper, a more useful unit of
design than a new architecture.

\subsection{Take-home sentence}

\begin{quote}
ST-JEWM does not prove that spike traces are the highest-scoring
control representation. It proves that post-spike traces can be
valid, calibrated, event-aligned predictive states under a stricter
membrane-forbidden world-model interface, and that such claims require
collapse-robust diagnostics rather than latent cosine success alone.
\end{quote}

% ============================================================================
\section*{Claim Control Table}

\begin{table}[ht]
\centering
\small
\begin{tabular}{llp{8cm}}
\toprule
Claim & Status & Evidence \\
\midrule
STJEWM competitive on env-native success & supported & 67.1 vs best 69.5 (specialist AVG); $\pm 4$pp on G16 generalist \\
STJEWM raw SOTA & \textbf{not claimed} & never the best on any single suite \\
LeWM-SR can be inflated by collapse & supported & MLP LeWM-SR 95.5\%, div 0.0002 $\to$ collapse signature \\
STJEWM event-aligned latents (specialist + generalist) & supported & $\rho \geq 0.62$ specialist, $\rho \geq 0.99$ generalist (G4 / G8 / G16) \\
STJEWM calibrated, non-collapsed, non-noisy, non-over-reactive & supported & div $0.011 \pm 0.001$, resp $0.21 \pm 0.01$, $\rho \geq 0.99$ across all 6 readouts and 3 scales \\
Membrane-readout catastrophically fails stress & \textbf{refuted} in v0.7.2 & stress env-SR 25.5\% AVG, within 0.5pp of trace-only \\
Planner causally relies on event-window trace component & \textbf{refuted} & zeroing trace at event-aligned steps does not reduce env-SR more than random zeroing \\
Multi-seed std bars on generalist eval & deferred & 1-seed numbers reported honestly \\
\bottomrule
\end{tabular}
\end{table}

% ============================================================================
\appendix

\section*{Appendix outline}

\textbf{Appendix A --- Implementation details.} Environment list (16
standard + 4 stress + LeWM), training hyperparameters, CEM planner
settings, all six readout-mode implementations, data generation
pipeline (multi-env spec format).

\textbf{Appendix B --- Specialist per-env tables.} All 20 standard
envs $\times$ 13 models per metric (env-SR, LeWM-SR); all 4 stress
envs.

\textbf{Appendix C --- Generalist per-suite tables.} G4 / G8 / G16
raw matrices for env-SR, LeWM-SR, divergence, responsiveness,
event-align $\rho$ per (env, model) cell.

\textbf{Appendix D --- Event probe details.} Event labels, AUROC per
env-target, probe training protocol, hold-out splits.

\textbf{Appendix E --- Event alignment details.} Event-boundary
definition (observation first-difference local maxima), Pearson
$\rho$ per env, computation protocol.

\textbf{Appendix F --- Collapse-diagnostic derivation.} Toy proof that
a constant latent has div = 0 regardless of planner. Responsiveness
definition with bounded-gain caveat. Alignment-vs-divergence
independence argument.

\textbf{Appendix G --- Historical experiments and claim revision
audit.} v0.4 ``membrane collapses under stress'' $\to$ v0.7.2
refutation. v0.7.4 ``LeWM-SR collapse signature invisible'' $\to$
v0.7.5 diagnostic. v0.7.3 pilot (``4-env subset too small'')
$\to$ v0.7.5 16-env generalist. This appendix makes the empirical
story self-auditing rather than history-of-projects.

% ============================================================================
\section*{References}

\small
[1] LeWM-style joint-embedding world model (Hafner et al., 2023; this paper's repo, \texttt{README.md}).

[2] CuBiFAE (Kaiser et al., 2024 ICML).

[3] SpikeDreamer (Hong et al., 2024 AAAI).

[4] SLT-LIF (Liu et al., 2024 NeurIPS workshop).

[5] STJEWM \texttt{ReadoutMode} enumeration, \texttt{code/stjewm.py} line 50 (\texttt{RATE\_ONLY} docstring: ``moving\_average(spike).detach()'').

\end{document}
